% Options for packages loaded elsewhere
\PassOptionsToPackage{unicode}{hyperref}
\PassOptionsToPackage{hyphens}{url}
\documentclass[
  ignorenonframetext,
]{beamer}
\newif\ifbibliography
\usepackage{pgfpages}
\setbeamertemplate{caption}[numbered]
\setbeamertemplate{caption label separator}{: }
\setbeamercolor{caption name}{fg=normal text.fg}
\beamertemplatenavigationsymbolsempty
% remove section numbering
\setbeamertemplate{part page}{
  \centering
  \begin{beamercolorbox}[sep=16pt,center]{part title}
    \usebeamerfont{part title}\insertpart\par
  \end{beamercolorbox}
}
\setbeamertemplate{section page}{
  \centering
  \begin{beamercolorbox}[sep=12pt,center]{section title}
    \usebeamerfont{section title}\insertsection\par
  \end{beamercolorbox}
}
\setbeamertemplate{subsection page}{
  \centering
  \begin{beamercolorbox}[sep=8pt,center]{subsection title}
    \usebeamerfont{subsection title}\insertsubsection\par
  \end{beamercolorbox}
}
% Prevent slide breaks in the middle of a paragraph
\widowpenalties 1 10000
\raggedbottom
\AtBeginPart{
  \frame{\partpage}
}
\AtBeginSection{
  \ifbibliography
  \else
    \frame{\sectionpage}
  \fi
}
\AtBeginSubsection{
  \frame{\subsectionpage}
}
\usepackage{iftex}
\ifPDFTeX
  \usepackage[T1]{fontenc}
  \usepackage[utf8]{inputenc}
  \usepackage{textcomp} % provide euro and other symbols
\else % if luatex or xetex
  \usepackage{unicode-math} % this also loads fontspec
  \defaultfontfeatures{Scale=MatchLowercase}
  \defaultfontfeatures[\rmfamily]{Ligatures=TeX,Scale=1}
\fi
\usepackage{lmodern}
\ifPDFTeX\else
  % xetex/luatex font selection
\fi
% Use upquote if available, for straight quotes in verbatim environments
\IfFileExists{upquote.sty}{\usepackage{upquote}}{}
\IfFileExists{microtype.sty}{% use microtype if available
  \usepackage[]{microtype}
  \UseMicrotypeSet[protrusion]{basicmath} % disable protrusion for tt fonts
}{}
\makeatletter
\@ifundefined{KOMAClassName}{% if non-KOMA class
  \IfFileExists{parskip.sty}{%
    \usepackage{parskip}
  }{% else
    \setlength{\parindent}{0pt}
    \setlength{\parskip}{6pt plus 2pt minus 1pt}}
}{% if KOMA class
  \KOMAoptions{parskip=half}}
\makeatother
\usepackage{longtable,booktabs,array}
\newcounter{none} % for unnumbered tables
\usepackage{calc} % for calculating minipage widths
\usepackage{caption}
% Make caption package work with longtable
\makeatletter
\def\fnum@table{\tablename~\thetable}
\makeatother
\usepackage{graphicx}
\makeatletter
\newsavebox\pandoc@box
\newcommand*\pandocbounded[1]{% scales image to fit in text height/width
  \sbox\pandoc@box{#1}%
  \Gscale@div\@tempa{\textheight}{\dimexpr\ht\pandoc@box+\dp\pandoc@box\relax}%
  \Gscale@div\@tempb{\linewidth}{\wd\pandoc@box}%
  \ifdim\@tempb\p@<\@tempa\p@\let\@tempa\@tempb\fi% select the smaller of both
  \ifdim\@tempa\p@<\p@\scalebox{\@tempa}{\usebox\pandoc@box}%
  \else\usebox{\pandoc@box}%
  \fi%
}
% Set default figure placement to htbp
\def\fps@figure{htbp}
\makeatother
\setlength{\emergencystretch}{3em} % prevent overfull lines
\providecommand{\tightlist}{%
  \setlength{\itemsep}{0pt}\setlength{\parskip}{0pt}}
% Slide layout defaults for warning-clean scientific decks.
\usepackage{etoolbox}
\IfFileExists{xurl.sty}{\usepackage{xurl}}{}
\IfFileExists{seqsplit.sty}{\usepackage{seqsplit}}{\newcommand{\seqsplit}[1]{#1}}
\protected\def\breakseq#1{\seqsplit{#1}}
\protected\def\breaktt#1{\begingroup\ttfamily\seqsplit{#1}\endgroup}
\setlength{\emergencystretch}{6em}
\tolerance=5000
\hbadness=10000
\hfuzz=1pt
\setlength{\tabcolsep}{2pt}
\AtBeginEnvironment{longtable}{\tiny\renewcommand{\arraystretch}{0.86}\setlength{\tabcolsep}{1pt}}
\AtBeginEnvironment{tabular}{\tiny\renewcommand{\arraystretch}{0.86}\setlength{\tabcolsep}{1pt}}

% Natbib and cross-reference fallbacks — slides don't load natbib
% or cleveref, but combined-PDF manuscript prose may emit these
% commands. The fallback renders citations as a bracketed key list
% and cross-references as detokenized labels so slides don't fail on
% undefined control sequences or raw underscores. \providecommand is
% a no-op if packages load later.
\providecommand{\citep}[1]{[#1]}
\providecommand{\citet}[1]{#1}
\providecommand{\citealp}[1]{#1}
\providecommand{\citeauthor}[1]{#1}
\providecommand{\citeyear}[1]{#1}
\providecommand{\cref}[1]{\texttt{\detokenize{#1}}}
\providecommand{\Cref}[1]{\texttt{\detokenize{#1}}}

\newtheorem{proposition}{Proposition}
\newtheorem{hypothesis}{Hypothesis}
\usepackage{bookmark}
\IfFileExists{xurl.sty}{\usepackage{xurl}}{} % add URL line breaks if available
\urlstyle{same}
% fallback for those not using the hyperref driver hyperxmp:
\makeatletter
\@ifundefined{xmpquote}{\newcommand{\xmpquote}[1]{#1}}{}
\makeatother
\hypersetup{
  hidelinks,
  pdfcreator={LaTeX via pandoc}}

\author{\texorpdfstring{}{}}
\date{}

\begin{document}

\section{Results}\label{sec:results}

\begin{frame}[fragile,allowframebreaks]{Results}
Results are interpreted only after the registered analysis has run and
the deviation ledger ({[}Section @sec:deviations{]}) has classified any
departures from the plan. Every value below is read from
\breaktt{output/data/demo\_analysis.json}, which is regenerated by
\breaktt{scripts/generate\_figures.py} from the tested
\breaktt{run\_registered\_analysis} function.
\end{frame}

\begin{frame}[fragile,allowframebreaks]{Confirmatory outcome (H1,
\texttt{primary\_score})}
\protect\phantomsection\label{confirmatory-outcome-h1-primary_score}
On the deterministic demonstration data, the registered two-sided
permutation test returns:

{\def\LTcaptype{none} % do not increment counter
\begin{longtable}[]{@{}ll@{}}
\toprule\noalign{}
Quantity & Value \\
\midrule\noalign{}
\endhead
Control mean & \texttt{-0.178991} \\
Treatment mean & \texttt{0.823948} \\
Observed mean difference & \texttt{1.002940} \\
Permutations & \texttt{2000} \\
Shuffles with absolute difference ≥ observed & \texttt{0} \\
Two-sided permutation p-value & \texttt{0.0005} \\
Preregistered \texttt{alpha} & \texttt{0.05} \\
Confirmatory decision & significant \\
\bottomrule\noalign{}
\end{longtable}
}

The observed difference of \texttt{1.003} lies far in the right tail of
the permutation null distribution: none of the \texttt{2000} label
shuffles produced a mean difference as extreme, so the add-one-corrected
p-value is \texttt{0.0005}, below the preregistered
\texttt{alpha\ =\ 0.05}. H1 is therefore \textbf{supported} as a
confirmatory result. \breakseq{[@fig:permutation\_result]} shows the observed
statistic against the null it was tested against.

\begin{figure}
\centering
\pandocbounded{\includegraphics[keepaspectratio,alt={Registered permutation test on the deterministic demonstration data, rendered by plot\_permutation\_result from permutation\_null\_distribution and run\_permutation\_test. The histogram is the label-permutation null distribution of the treatment-minus-control mean difference (2000 seeded shuffles); the solid red line marks the observed difference of 1.003 and the dotted line its mirror image. The annotation reports n = 24 per group, 0 shuffles at least as extreme, a two-sided p-value of 0.0005, and alpha = 0.05. The observed statistic sits outside the entire null, which is why the primary confirmatory claim holds.}]{../figures/permutation_result.png}}
\caption{Registered permutation test on the deterministic demonstration
data, rendered by \protect\breaktt{plot\_permutation\_result} from
\protect\breaktt{permutation\_null\_distribution} and
\protect\breaktt{run\_permutation\_test}. The histogram is the
label-permutation null distribution of the treatment-minus-control mean
difference (\protect\texttt{2000} seeded shuffles); the solid red line
marks the observed difference of \protect\texttt{1.003} and the dotted
line its mirror image. The annotation reports \protect\texttt{n\ =\ 24}
per group, \protect\texttt{0} shuffles at least as extreme, a two-sided
p-value of \protect\texttt{0.0005}, and \protect\texttt{alpha\ =\ 0.05}.
The observed statistic sits outside the entire null, which is why the
primary confirmatory claim holds.}\label{fig:permutation_result}
\end{figure}
\end{frame}

\begin{frame}[fragile,allowframebreaks]{Claim boundary}
\protect\phantomsection\label{claim-boundary}
The registration preregisters exactly one confirmatory outcome. The
review packet built by \breaktt{build\_review\_packet} records
\texttt{primary\_score} as the sole confirmatory outcome and returns an
integrity score of \texttt{1.0} for a plan-faithful execution. Any other
endpoint examined during analysis is exploratory and is reported under
{[}Section @sec:deviations{]}, never here. This separation is enforced
in code, not merely asserted in prose.
\end{frame}

\end{document}
